%% Chapter-by-chapter skeleton of the thesis, one bullet per argument move.
%% Written on the CSM thesis template (HYthesisML.cls) of the Department of
%% Computer Science; see CSM-thesis-template_2021_09_07/ for the pristine copy.
%% Kept in step with Ha_Vien_Mastersthesis_2026.tex, which uses the same class and options.

% STEP 1: language, pagination, programme, and study track.
% openright is passed through to report so chapters start on a right-hand page;
% the template's own default is openany.
% a4paper has to be given here: the template sets it with a \PassOptionsToClass
% placed after \documentclass, by which time report is already loaded, so the
% stock template silently produces US Letter pages.
\documentclass[english,twoside,openright,a4paper,censored,csm,software-track-2020]{HYthesisML}

\usepackage{csquotes}

% REFERENCES: numeric (Vancouver) style, one of the three styles the template
% offers. sorting=none keeps the numbers in citation order, as IEEEtran did.
\usepackage[style=numeric,bibstyle=numeric,backend=biber,natbib=true,
            sorting=none,sortcites=true,maxbibnames=99,giveninits=true,uniquename=init]{biblatex}
\addbibresource{bibliography.bib}
% List the authors in the bibliography as Lastname F.
\DeclareNameAlias{sortname}{family-given}
% Without a mark of its own the bibliography keeps the last chapter's title in
% the running head of its continuation pages.
\defbibheading{bibliography}[\bibname]{%
  \chapter*{#1}%
  \markboth{#1}{#1}}

% lmodern, textcomp, color, graphicx, epsfig, subfigure, tikz, amsmath, amssymb,
% caption, titlesec and appendix are all loaded by the class itself.
\usepackage[pdftex, plainpages=false]{hyperref} % For hyperlinks and pdf metadata
\usepackage{fancyhdr} % For nicer page headers
%\usepackage{wrapfig} % For nice text-wrapping figures (use at own discretion)

\onehalfspacing %line spacing
%\singlespacing
%\doublespacing

%\fussy 
\sloppy % sloppy and fussy commands can be used to avoid overlong text lines

% STEP 2:
% Set up all the information for the title page and the abstract form.
% Replace parameters with your information.
\title{Performance, Safety, and Engineering Trade-offs of Rust and C++ in Nokia RAN Control-Plane Systems}
\author{Ha Vien}
\date{\today}
% Supervisors are printed on the abstract form rather than on the title page.
% TODO before submission: the template asks for the academic title (Prof., Dr.)
% in front of each name.
\supervisors{Keijo Heljanko, Christoph Luneburg, Petteri Kotimaa}
% The class prints examiners nowhere, so this is recorded but has no output.
% The faculty has not named the second examiner yet.
\examiners{Prof. Keijo Heljanko}
\keywords{Rust, C++, RAN control plane, tail latency, memory safety, fault injection, incremental adoption}
% Where deposited fills itself in from the class. Additional information is the
% study track, taken from the software-track-2020 class option.
\additionalinformation{\translate{\track}}


\classification{\protect{\ \\
\  Software and its engineering $\rightarrow$ Software notations and tools $\rightarrow$ General programming languages\  \\
\  Computer systems organization $\rightarrow$ Real-time systems $\rightarrow$ Real-time system specification\\
\  Networks $\rightarrow$ Network types $\rightarrow$ Mobile networks
}}

% if you want to quote someone special. You can comment this line and there will be nothing on the document.
%\quoting{Bachelor's degrees make pretty good placemats if you get them laminated.}{Jeph Jacques} 


% OPTIONAL STEP: Set up properties and metadata for the pdf file that pdfLaTeX makes.
% But you don't really need to do this unless you want to.
\hypersetup{
    %bookmarks=true,         % show bookmarks bar first?
    unicode=true,           % to show non-Latin characters in Acrobat’s bookmarks
    pdftoolbar=true,        % show Acrobat’s toolbar?
    pdfmenubar=true,        % show Acrobat’s menu?
    pdffitwindow=false,     % window fit to page when opened
    pdfstartview={FitH},    % fits the width of the page to the window
    pdftitle={Performance, Safety, and Engineering Trade-offs of Rust and C++ in Nokia RAN Control-Plane Systems},
    pdfauthor={Ha Vien},
    pdfsubject={Master's thesis, Computer Science},
    pdfcreator={Ha Vien},
    pdfproducer={pdfLaTeX}, % producer of the document
    pdfkeywords={Rust} {C++} {RAN control plane} {tail latency} {memory safety},
    pdfnewwindow=true,      % links in new window
    colorlinks=true,        % false: boxed links; true: colored links
    linkcolor=black,        % color of internal links
    citecolor=black,        % color of links to bibliography
    filecolor=magenta,      % color of file links
    urlcolor=cyan           % color of external links
}

% The same three generated macro files Ha_Vien_Mastersthesis_2026.tex reads: campaign
% summary numbers from bench/plot_figures.py, the size and complexity proxies from
% bench/complexity.py, and the artefact-level measurements from bench/artefact_facts.py.
% The results bullets quote these rather than carrying their own copies of the numbers,
% so a rerun of bench/refresh_thesis.sh updates the skeleton with the thesis.
% Generated by bench/plot_figures.py -- do not edit by hand.
\newcommand{\ReprLoad}{28}
\newcommand{\ReprLoadAch}{28.0}
\newcommand{\ReprLoadShortfall}{0.1}
\newcommand{\SatKnee}{334}
\newcommand{\TrackingLoad}{336}
\newcommand{\NRepeats}{20}
\newcommand{\NLoadLevels}{14}
\newcommand{\NVariants}{4}
\newcommand{\NRuns}{1120}
\newcommand{\PNineNineSpread}{14}
\newcommand{\PFiftySpread}{2.9}
\newcommand{\PFiftyRange}{3}
\newcommand{\PNineNineRange}{460}
\newcommand{\PNineNineMaxCI}{2540}
\newcommand{\MaxOfferedLoad}{784}
\newcommand{\NTests}{10}
\newcommand{\MinAttainableP}{0.0000019}
\newcommand{\CorrectedFloor}{0.000019}
\newcommand{\NSignificant}{2}
\newcommand{\SmallestHolm}{0.00010}
\newcommand{\NExclZero}{4}
\newcommand{\NTailExclZero}{1}
\newcommand{\TailExclZeroRows}{p99 latency under system}
\newcommand{\RobustExclZeroRows}{p50 latency, Peak RSS and CPU utilisation under system}
\newcommand{\PFiftyDiffAbsMax}{1.6}
\newcommand{\PFiftyDiffCIMax}{3}
\newcommand{\RssLangDiffLo}{1.3}
\newcommand{\RssLangDiffHi}{1.7}
\newcommand{\RssLangCILo}{-0.5}
\newcommand{\RssLangCIHi}{1.9}
\newcommand{\CpuDiffAbsMax}{3.4}
\newcommand{\CrossAllocConsistent}{p50 latency, p99 latency, p99.9 latency, Peak RSS and CPU utilisation}
\newcommand{\CrossAllocReversed}{none}
\newcommand{\NCrossAllocReversed}{0}
\newcommand{\RssAllocDiffLo}{55.9}
\newcommand{\RssAllocDiffHi}{56.5}
\newcommand{\BudgetRepeats}{15}
\newcommand{\BudgetCoreCpp}{34.4}
\newcommand{\BudgetCoreRust}{21.3}
\newcommand{\BudgetSvcCpp}{13.52}
\newcommand{\BudgetSvcRust}{12.63}
\newcommand{\BudgetHandoffCpp}{13.49}
\newcommand{\BudgetHandoffRust}{12.61}
\newcommand{\BudgetCoreShareMax}{0.250}
\newcommand{\BudgetCoreFaster}{Rust}
\newcommand{\BudgetCoreRatio}{1.6}
\newcommand{\BudgetHandoffSlower}{C++}
\newcommand{\BudgetBoundVsCore}{101}
\newcommand{\FfiCrossing}{2.54}
\newcommand{\FfiOpRepeats}{15}
\newcommand{\FfiOpNative}{12.3}
\newcommand{\FfiOpSeam}{12.7}
\newcommand{\FfiOpDelta}{+192}
\newcommand{\FfiOpDeltaLo}{-954}
\newcommand{\FfiOpDeltaHi}{+1341}
\newcommand{\FfiOpP}{0.239}
\newcommand{\FfiOpBoundPct}{10.91}

% Generated by bench/complexity.py -- do not edit by hand.
\newcommand{\CppFuncs}{37}
\newcommand{\RustFuncs}{40}
\newcommand{\CppNloc}{511}
\newcommand{\RustNloc}{608}
\newcommand{\CppCcnTotal}{93}
\newcommand{\RustCcnTotal}{79}
\newcommand{\CppCcnMean}{2.5}
\newcommand{\RustCcnMean}{2.0}
\newcommand{\CppCcnMax}{13}
\newcommand{\RustCcnMax}{6}
\newcommand{\CppCcnMaxFn}{\texttt{worker\_loop}}
\newcommand{\RustCcnMaxFn}{\texttt{add\_drb}}
\newcommand{\CcnTotalDrop}{15}
\newcommand{\NlocRise}{19}

% Generated by bench/artefact_facts.py -- do not edit by hand.
\newcommand{\CppUeCtxBytes}{96}
\newcommand{\RustUeCtxBytes}{88}
\newcommand{\UeCtxGapBytes}{8}
\newcommand{\CppUeCtxStatePad}{3}
\newcommand{\CppUeCtxKeyEnd}{50}
\newcommand{\CppUeCtxKeyPad}{6}
\newcommand{\CppRequestBytes}{160}
\newcommand{\RustRequestBytes}{96}
\newcommand{\SignalingMsgBytes}{68}
\newcommand{\SweepUes}{5\,000}
\newcommand{\UeCtxGapTotalKB}{40}
\newcommand{\CppQueueResidentMB}{1.3}
\newcommand{\RustQueueResidentMB}{0.8}
\newcommand{\CppLineCov}{94.4}
\newcommand{\RustLineCov}{97.9}
\newcommand{\RustRegionCov}{99.1}
\newcommand{\ConfHandOps}{223}
\newcommand{\ConfGenOps}{9\,000}
\newcommand{\FfiSeamLines}{231}
\newcommand{\FfiHeaderLines}{67}
\newcommand{\FfiEntryPoints}{11}
\newcommand{\FfiUnsafeFns}{9}
\newcommand{\DeadlockTsanHits}{3}
\newcommand{\DeadlockTsanRuns}{12}
\newcommand{\FaultClasses}{14}
\newcommand{\FaultClassesPrevented}{7}
\newcommand{\FaultClassesBoundary}{7}
\newcommand{\FaultClassesUndetected}{4}
\newcommand{\LostUpdateExpected}{40\,000}
\newcommand{\LostUpdateRustHandled}{28\,714}
\newcommand{\LostUpdateRustLost}{11\,286}
\newcommand{\LostUpdateCppHandled}{27\,320}
\newcommand{\LostUpdateCppLost}{12\,680}
\newcommand{\CppTestBinaries}{6}
\newcommand{\RustTests}{27}

% Generated by bench/adoption_cost.py -- do not edit by hand.
\newcommand{\AdoptRepeats}{15}
\newcommand{\AdoptJobs}{14}
\newcommand{\BuildLibCppWall}{1.9}
\newcommand{\BuildLibRustWall}{1.2}
\newcommand{\BuildLibCppCpu}{2.6}
\newcommand{\BuildLibRustCpu}{3.9}
\newcommand{\BuildLibWallRatio}{1.7}
\newcommand{\BuildLibCpuRatio}{1.5}
\newcommand{\BuildLibWallFaster}{Rust}
\newcommand{\BuildLibCpuCheaper}{C++}
\newcommand{\BuildAllCppWall}{5.7}
\newcommand{\BuildAllRustWall}{8.1}
\newcommand{\BuildAllCppCpu}{35.9}
\newcommand{\BuildAllRustCpu}{64.8}
\newcommand{\BuildAllCpuRatio}{1.8}
\newcommand{\BuildAllFirstPartyCpu}{29.6}
\newcommand{\BuildAllFirstPartyWall}{3.8}
\newcommand{\BuildDepShareCpu}{54}
\newcommand{\BuildAllFirstPartyCpuRatio}{0.8}
\newcommand{\BuildIncCppTu}{1.2}
\newcommand{\BuildIncCppHeader}{5.8}
\newcommand{\BuildIncRust}{3.6}
\newcommand{\BuildIncTuRatio}{3.1}
\newcommand{\BuildIncHeaderRatio}{1.6}
\newcommand{\BuildConfigure}{0.91}
\newcommand{\FootTarget}{conformance\_runner}
\newcommand{\FootCppShipped}{106}
\newcommand{\FootRustShipped}{723}
\newcommand{\FootShippedRatio}{6.8}
\newcommand{\FootCppStatic}{1577}
\newcommand{\FootRustStatic}{723}
\newcommand{\FootStaticRatio}{0.5}
\newcommand{\FootCppDynamic}{106}
\newcommand{\FootRustDynamic}{269}
\newcommand{\FootDynamicRatio}{2.5}
\newcommand{\FootRustStdShare}{63}
\newcommand{\FootCppStdShare}{93}
\newcommand{\FootSharedCpp}{2531}
\newcommand{\FootSharedRust}{5295}
\newcommand{\FootStaticSmaller}{Rust}
\newcommand{\FootDynamicSmaller}{C++}
\newcommand{\FootStaticRatioAbs}{2.2}
\newcommand{\FootCppLocaleShare}{47}
\newcommand{\DepsCore}{0}
\newcommand{\DepsJemalloc}{8}
\newcommand{\DepsHarness}{17}
\newcommand{\DepsUnsafeJemalloc}{623}
\newcommand{\DepsUnsafeHarness}{1064}
\newcommand{\CoreUnsafe}{0}
\newcommand{\DepsJemallocShipping}{5}
\newcommand{\DepsJemallocBuildOnly}{3}
\newcommand{\DepsHarnessShipping}{16}
\newcommand{\DepsHarnessBuildOnly}{1}


\begin{document}

% Generate title page.
\maketitle

% STEP 3:
% Write your abstract (of course you really do this last).
% You can make several abstract pages (if you want it in different languages)
% by repeating the otherlanguage environment with a different language.

% The abstract form is a fixed-height box, so it is set single-spaced as in the
% template; the 1.5 spacing of the body would cost it a fifth of its lines.
{\singlespacing
\begin{otherlanguage}{english}
\begin{abstract}
\noindent Abstract to be written last.
\end{abstract}
\end{otherlanguage}
}

% Place ToC
\mytableofcontents

\mainmatter

% -----------------------------------------------------------------------------------
% STEP 4: Write the thesis.
% Your actual text starts here. You shouldn't mess with the code above the line except
% to change the parameters. Removing the abstract and ToC commands will mess up stuff.
\chapter{Introduction}

\begin{itemize}
  \item The 5G RAN control plane (the gNB-CU-CP) now runs as software on ordinary servers.
        It is large, it keeps state for every connected device, it serves many requests at
        once, and the 5G target for control plane latency, the transition from a
        battery-efficient state to continuous data transfer, is 10~ms~\cite{3gpp38913}.
        One defect, or one regression in the slowest requests, therefore degrades the whole
        site. The work is carried out at Nokia on a UE Context Management service modelled on a
        real gNB-CU-CP component.
  \item This software is written mostly in C++. C++ gives the control over memory and timing
        that such a system needs, but it leaves memory-safety and concurrency mistakes to the
        programmer~\cite{szekeres2013sok}. Rust promises the same performance while rejecting
        those mistakes at compile time~\cite{klabnik2023rust}. The choice has a long horizon:
        3GPP has allocated Release 20 to the 6G studies and Release 21 to the normative work,
        with migration from 5G in scope~\cite{3gpp2024road6g}, so a language chosen for the
        control plane of the current generation is inherited by the next.
  \item The problem is a lack of evidence. Telecom vendors are under growing pressure to move
        to memory-safe languages~\cite{cisa2023memorysafe}, and the closest prior work is a
        cluster of at least five Ericsson master's theses, none of which answers the question
        as posed here. Three keep the runtime in play alongside the language, comparing
        Boost.Fibers, goroutines, Tokio, Monoio, Glommio and DPDK Eventdev rather than the
        languages that host
        them~\cite{karlback2021ran,olsson2023asyncpackets,arvidsson2025async}. The other two
        report execution time and peak memory rather than tail behaviour under rising
        load~\cite{olsson2023ranmgmt,emberg2023asn1}. Emberg and Wadsten come closest in
        method, with a matched C++ and Rust pair joined by a C wrapper, but measure a
        stateless transformation~\cite{emberg2023asn1}. No study holds the specification,
        concurrency model, and build settings fixed for a stateful, concurrent subsystem while
        sweeping past saturation and measuring incremental migration cost.
  \item That evidence is also concentrated in one vendor and untested outside it, while on the
        other side there is no study of Rust among Nokia's indexed publications at all. The
        nearest peer-reviewed position from that side argues for making memory-error sanitisation
        practical rather than for replacing it with a compile-time
        guarantee~\cite{gjomemo2016mesc}, which is the strategy this thesis tests against.
  \item Four research questions frame the work. RQ1 compares latency predictability (median,
        p99, p99.9, maximum) under increasing load and overload with the concurrency model held
        fixed~\cite{dean2013tail}. RQ2 asks which defect classes Rust removes at compile time,
        at what stage the rest are detected, and which classes it does not prevent. RQ3 asks what
        overhead Rust's safety model imposes on throughput, CPU, and memory. RQ4 asks what
        incremental adoption costs in effort, FFI interoperability, and risk inside an existing
        C++ RAN codebase.
  \item The work contributes five things. A reproducible, telecom-relevant benchmarking
        framework. A frozen shared specification with two parity-tested implementations. A
        quantitative comparison that emphasises tail latency under overload. A defect-class
        taxonomy that marks where Rust's compile-time guarantee stops. And evidence-based
        guidelines for incremental adoption, where the dominant cost is the learning curve and not
        run-time~\cite{fulton2021adopting}.
\end{itemize}

\chapter{Background: Systems Languages for Latency-Critical Software\label{chapter:background}}
% The class prints the chapter titles in the running head in capitals, where the
% full title of this chapter runs into the margin.
\chaptermark{Background}

\begin{itemize}
  \item In the NG-RAN architecture the base station (gNB) is split into a central and a
        distributed unit, and the central unit is split again into a control plane and a user
        plane. UE Context Management sits in the control plane, the gNB-CU-CP, and reaches its
        neighbours over the standardised F1, E1, and NG interfaces~\cite{3gpp38401}.
  \item Deadlines are met or missed because of predictability, not average speed. The metrics
        that matter are therefore the p99 and p99.9 latencies and the jitter around
        them~\cite{dean2013tail}. They have to be measured past the saturation point and
        recorded without coordinated omission, that is, without letting a load generator that
        has fallen behind hide its own slowest
        requests~\cite{schroeder2006openclosed,tene_hdrhistogram}.
  \item C++ offers abstractions that cost nothing at run time, deterministic cleanup through
        RAII, and no garbage collector. Safety, however, remains the programmer's
        responsibility~\cite{szekeres2013sok}, and sanitizers report only the errors that a
        given test run happens to reach.
  \item Rust has the same run-time model and also no garbage collector, but it checks memory
        safety and freedom from data races at compile time, using ownership, borrowing, and
        thread-safety traits such as \texttt{Send}~\cite{klabnik2023rust,jung2021safe}. Code
        that needs to step outside
        those rules must be marked \texttt{unsafe}, which keeps such code small and easy to
        find.
  \item The comparison is therefore C++ against Rust. Garbage-collected languages are left out
        because their pauses are incompatible with the latency targets. The concurrency model
        is held fixed in both implementations, and the C ABI is treated as the realistic
        adoption path. No earlier study isolates the language this way for a stateful,
        concurrent subsystem~\cite{karlback2021ran,emberg2023asn1}.
\end{itemize}

\chapter{Problem Context: UE Context Management and Predictability\label{chapter:context}}
\chaptermark{Problem Context}

\begin{itemize}
  \item UE Context Management owns the lifecycle of one context per connected device. Each
        context moves through four states, Idle, Connected, Reconfiguring, and Released, driven
        by a small set of lifecycle operations. An operation that is illegal in the current
        state returns a typed error instead of being carried out.
  \item The subsystem was chosen because it is both representative and small enough to keep
        under control. It has mutable state per device, a strict lifecycle, heavy concurrency,
        and hard latency targets~\cite{3gpp38913,3gpp22261}, and it is still simple enough to specify
        precisely and to test for identical behaviour in two languages.
  \item Five defect classes fall inside Rust's guarantee, and each has a natural form in this
        subsystem: use-after-free, double-free, iterator invalidation, data race, and
        unguarded shared mutation. What the guarantee rests on is aliasing control and
        data-race freedom~\cite{jung2021safe}. The catalogue is this study's own.
  \item Five further classes are included on purpose from outside that guarantee: a lock-order
        deadlock, a lost update, a leak through a reference cycle, an integer overflow in a
        release build, and the blast radius of a failed shard worker. They keep the study from
        confirming a foregone conclusion and show instead where the guarantee stops. Two of
        them return the opposite of the expected result, with the C++ tooling reporting what
        Rust's does not.
  \item Scale is what makes the tail matter. One procedure triggers many operations and one
        site serves many devices, so a p99.9 measured on a single operation becomes a steady
        rate of degraded procedures in the field~\cite{dean2013tail}. This is why the load
        sweep is pushed past the point where the service starts to fall behind.
\end{itemize}

\chapter{Objectives\label{chapter:objectives}}

\begin{itemize}
  \item A documented literature search of the citation databases and of the Nordic thesis
        repositories, where the closest work turns out to live, found no study that holds a
        stateful control-plane subsystem fixed while measuring tail latency of Rust and C++
        past saturation under the same concurrency
        model~\cite{karlback2021ran,arvidsson2025async,emberg2023asn1}.
  \item Five objectives follow from the research questions: measure predictability (O1),
        establish which defect classes are prevented (O2), measure the resource overhead (O3),
        build and measure a migration prototype (O4), and publish reusable artifacts and
        guidance (O5).
  \item Three hypotheses are stated in advance, so that the measurements can refute them. H1:
        below saturation the two implementations perform the same. The prior evidence points to
        near-parity with a measurable residual rather than to equality, a Rust network driver
        running 4\% slower than its C reference while executing 63\% more instructions per
        packet~\cite{emmerich2019drivers}.
        H2: Rust prevents at compile time the defect classes for which C++ can only offer
        best-effort detection. H3: the cost of calling Rust from C++ is negligible.
  \item The work succeeds if the evidence is good enough to base a decision on: repeated runs
        reported with confidence intervals~\cite{georges2007rigorous}, defect claims shown on
        programs that really compile and run, migration cost read off a working prototype, and
        every number reproducible from the published artifacts.
  \item The scope is drawn narrowly to protect internal validity. The full gNB, the PHY and L2
        layers, the user plane, captured traffic traces, managed languages, and production
        operations are all left out.
\end{itemize}

\chapter{Research Methodology\label{chapter:methodology}}

\begin{itemize}
  \item The approach is empirical and constructive. Instead of comparing two systems that
        already exist and differ in many ways, the study builds a matched pair of
        implementations from one frozen specification, so that the language is the only
        systematic difference between them. Both use an identical thread-per-shard concurrency
        model, because that structuring itself changes tail
        latency~\cite{enberg2019threadpercore}.
  \item RQ1 and RQ3 are answered with an open-loop load sweep that runs from light load into
        overload on pinned cores. Open-loop means the generator keeps to its send schedule
        instead of waiting for replies. Latency is recorded in identically configured HDR
        histograms~\cite{tene_hdrhistogram}, and resource use is sampled through the operating
        system. The results so far come from a dedicated dual-socket Xeon host running as a
        KVM guest. That host resolves the overload divergence and the resource comparison, but
        not the far tail below saturation: the frequency governor is not observable from inside
        the guest and CPU steal is non-zero. Locking frequency requires the bare-metal
        campaign.
  \item RQ2 is answered by controlled fault injection~\cite{natella2016fault}. A fixed
        catalogue of defects is written in both languages, and the expected outcome of each one
        is recorded in advance, so that a surprise shows up as a discrepancy instead of being
        explained away afterwards.
  \item RQ4 is answered by a working prototype in which C++ drives the Rust core through a
        foreign-function interface, the seam being where cross-language memory-management
        defects concentrate~\cite{li2022crosslang}. It yields three measurements: how large the
        seam is, how much \texttt{unsafe} code it needs, and what one call across it costs.
  \item Fairness is the central experimental control: the specification is frozen, allocator
        and optimisation settings are matched, and every point where the two languages had to
        differ is written down in a forced-idiom log. Reporting follows the guidance that a
        single figure invites conclusions run-to-run variation does not
        support~\cite{georges2007rigorous}: medians with nonparametric bootstrap confidence
        intervals~\cite{efron1993bootstrap}, and exact paired permutation tests corrected for
        multiple comparisons by the Holm procedure~\cite{holm1979simple}. All data is
        synthetic, so no data-protection concerns arise.
\end{itemize}

\chapter{Design and Implementation\label{chapter:design}}

\begin{itemize}
  \item Nothing is implemented until the specification is frozen. It fixes the public API, the
        UE state machine, the data structures and their algorithmic strategy, the locking and
        sharding granularity, and the concurrency model. Both implementations conform to it, so
        what the comparison measures is the language and not a difference of design or of skill.
        The specification is itself a deliverable and the study's primary control against
        implementation bias.
  \item One concurrency model is used on both sides, shared-nothing and thread-per-shard,
        because that structuring is known to change tail latency by
        itself~\cite{enberg2019threadpercore}. The identifier space is partitioned into shards,
        each owned by a single worker thread that alone reads or mutates it, so no lock guards
        context state. Callers submit a request to the owning shard's bounded queue and wait on
        a single-use completion handle. The queue blocks when it is full, which is deliberate:
        a drop policy would hide overload, and this study needs overload to appear as latency.
        The only shared mutable state either version permits is one atomic counter, and no
        control flow depends on its value.
  \item Fairness is enforced below the design as well as in it. Both builds link the same
        \texttt{jemalloc} and keep a system-allocator build as a control, both compile at
        matched optimisation, and both record latency through one HDR histogram algorithm
        configured identically~\cite{tene_hdrhistogram}. Hashing is pinned on both sides,
        SplitMix64 for shard assignment and FNV-1a inside the map, because C++'s default hash
        of an integer is typically the identity function while Rust's is cryptographically
        strong. Left as they came, the two standard libraries would have been the measurement.
  \item Where the languages cannot satisfy a requirement the same way, the divergence is
        written into a forced-idiom log and then measured. Rust takes its bounded blocking
        queue from \texttt{sync\_channel}, while C++ through C++23 has no such primitive and
        assembles one from a mutex and two condition variables. C++ has no payload-carrying
        enumeration in this idiom either, so a request becomes a tagged struct holding every
        argument field any operation might need, $\CppRequestBytes$~bytes against Rust's
        $\RustRequestBytes$, answered through a \texttt{variant} the caller has to unwrap. The
        net of all of it is a Rust core $\NlocRise\%$ longer in code lines whose total
        cyclomatic complexity is $\CcnTotalDrop\%$ lower~\cite{lizard}, the branchiest function
        in either core being the C++ worker loop that dispatches that \texttt{variant}. The
        Rust library contains no \texttt{unsafe} at all.
  \item Parity is tested and not assumed. Language-neutral vectors run in both implementations
        and their transcripts are diffed operation by operation. Because a cross-diff cannot
        tell ``both correct'' from ``both wrong in the same way'', a third implementation in
        Python acts as an independent oracle: it has to reproduce all $\ConfHandOps$
        hand-derived expectations before it may generate any of its own, and a campaign of
        roughly $\ConfGenOps$ generated operations found no divergence. The transition matrix
        is enumerated exhaustively instead of sampled, the C++ side runs under AddressSanitizer,
        ThreadSanitizer and UndefinedBehaviorSanitizer and the Rust side under Miri, and line
        coverage is reported, so a clean sanitiser result is bounded by the code that actually
        ran.
\end{itemize}

\chapter{Experimental Setup and Evaluation\label{chapter:eval}}

\begin{itemize}
  \item The workload is synthetic and spec-justified. A deterministic driver walks each UE
        through the connection lifecycle the 3GPP specifications define, and because its draw
        order is pinned, the per-operation-kind counts come out identical in both languages for
        a given seed, which doubles as a parity signal. Each sweep point is repeated
        \NRepeats{} times with the four language and allocator variants interleaved inside every
        repeat, so that thermal or background drift falls on all of them alike, and every point
        is a median with a bootstrap $95\%$ confidence
        interval~\cite{georges2007rigorous,efron1993bootstrap}. The representative operating
        point is not chosen by eye but taken as the busiest swept load at which all four
        variants still hold p99 below $5$~ms, $\ReprLoad$~kops/s against a measured saturation
        knee of $\SatKnee$~kops/s. All of it runs on a dedicated dual-socket Xeon pinned to
        cores $2$--$15$, and that host bounds what the campaign may claim, because it is a KVM
        guest. The frequency governor and the turbo state are unobservable from inside it, so
        frequency stability is assumed instead of enforced, and CPU steal is non-zero at roughly
        the order of the latencies under study. Run-level statistics survive that, which is why
        the overload and the resource comparisons are resolved here. The far tail below
        saturation does not, so it is reported as unresolved and not as a null, and the
        bare-metal protocol that would settle it is specified and published.
  \item Below saturation the four variants sit within $\PFiftyRange\,\mu$s of one another at the
        median, $\PFiftySpread\%$, and the paired test resolves a difference of
        $\PFiftyDiffAbsMax\,\mu$s in Rust's favour under both allocators. Higher up the
        distribution the instrument runs out: the widest interval on a single variant's own p99
        spans $\PNineNineMaxCI\,\mu$s against $\PNineNineRange\,\mu$s of spread between all
        four, and where the uncertainty about one variant exceeds the gap between them, an
        ordering would be reporting the scheduler. A second limit is structural and no host
        would fix it. An end-to-end operation is $\BudgetCoreShareMax\%$ the work the language
        performs, the rest being a cross-thread handoff charged to the operating
        system~\cite{zhou2025oltp}, so this sweep cannot bound a language effect in either
        direction. Measured on its own the bare operation is $\BudgetCoreRatio\times$ faster in
        \BudgetCoreFaster{}, and the likely cause is the standard hash map's layout and nothing
        about safety.
  \item Past saturation the languages are not equivalent, and the margin is not marginal. At
        $\OverloadLoad$~kops/s offered, \OverloadSustainedFaster{} sustains
        $\OverloadThrCpp$~kops/s where \OverloadSustainedSlower{} sustains
        $\OverloadThrRust$, a factor of $\OverloadThrRatio$, in \OverloadThrAgree{} repeats
        under both allocators and with the Holm-corrected $p$ at the floor of the whole test
        family. The slower variant is not simply doing less work: it draws
        $\OverloadCpuPerOpRatio$ times the CPU per completed operation, which is the signature
        of effort spent without completing anything. Isolating the bounded queue and the
        completion handle on their own does not reproduce that ordering at significance, so the
        mechanism sits in the wider handoff path. This is the largest and most reproducible
        language effect in the study, and a comparison that stopped at nominal load would have
        concluded the opposite.
  \item Safety is measured by injecting ten defects, five from inside the region Rust claims to
        govern and five chosen because safe Rust is expected \emph{not} to prevent
        them~\cite{natella2016fault}. Every one of the first five is a compile error, so it
        cannot be written at all, where in C++ each is valid code whose defect appears only if a
        test drives it and no single sanitiser covers the five. All five boundary cases compile
        in safe Rust and misbehave, and on two of them the C++ tooling is the one that reports:
        LeakSanitizer flags the reference cycle that Rust leaks in silence, and
        UndefinedBehaviorSanitizer flags the overflow that Rust defines to wrap.
        ThreadSanitizer catches the lock-order inversion on a minority of runs only, because it
        builds its lock-order graph from acquisitions that succeed and a run whose first
        inversion blocks never closes the cycle. Two Rust model checkers do report both original
        boundary defects, but only after each was re-imported against the checker's own
        primitives and shrunk to model size~\cite{bornholt2021shardstore}, so that half of the
        boundary is caught at test time and after a rewrite.
  \item Resource use separates the variants where latency could not, and it holds two effects of
        unequal size. The Rust build costs $\RssLangDiffLo$--$\RssLangDiffHi$~MiB more peak
        resident set than the C++ build under either allocator, while linking \texttt{jemalloc}
        at all costs $\RssAllocDiffLo$--$\RssAllocDiffHi$~MiB in either language. Both are real
        and the allocator's is an order of magnitude the larger. The language difference is a
        fixed cost and not a per-context one, and no hardware counters were collected, so it is
        established without being explained. CPU inverts with the regime, marginally lower for
        Rust while the service copes and much higher per completed operation once it does not.
        The migration seam is small on every axis measured: $\FfiEntryPoints$ C entry
        points in $\FfiSeamLines$ lines of Rust, all raw-pointer work confined to
        $\FfiUnsafeFns$ functions with the core left free of \texttt{unsafe}, and no cost
        detected against the same operation called natively, which bounds any per-operation tax
        at $\FfiOpBoundPct\%$. Three adoption costs beyond the seam were measured for the first
        time in this pass, and two contradict the received objection: $\BuildDepShareCpu\%$ of the
        Rust clean build is third-party code, so the project's own source compiles in
        $\BuildAllFirstPartyCpuRatio$ times the CPU \texttt{g++} needs, and the
        $\FootShippedRatio\times$ binary-size gap is a linkage policy that reverses when the policy
        is held fixed. The compared component resolves $\DepsCore$ crates and contains
        $\CoreUnsafe$ \texttt{unsafe} sites.
\end{itemize}

\chapter{Discussion\label{chapter:discussion}}

\begin{itemize}
  \item RQ1's answer divides at saturation, and reporting either half on its own would
        misrepresent it. Below the knee the two languages are equivalent in central behaviour by
        a margin no deployment would notice, which is a result about services of this shape, and
        the far tail in the same regime stays unresolved, which is an admission about the
        instrument. Past the knee they diverge by $\OverloadThrRatio\times$. The closest prior
        work measures this quantity for 5G control-plane procedures, finds the same shape, and
        traces it to stackless \texttt{async}/\texttt{await}~\cite{arvidsson2025async}. That
        attribution is accepted here. The point is that the shape appeared with no asynchronous
        runtime on either side and the concurrency model held fixed, so a runtime cannot be the
        only place such an effect originates. The prediction registered in advance also named
        where to look first, bounds and overflow checking, and that part of it was wrong.
        Recording the miss is worth more to the next study than quietly substituting the
        explanation that survived.
  \item On safety the difference is one of kind, and the boundary is reported as a first-class
        result and not as a caveat. A compiler that refuses to build a program rules the defect
        out of every execution, where a sanitiser rules it out only of the executions an
        instrumented test happens to take. What is absolute there is the execution coverage, not
        the breadth of the catalogue. Rust reasons about aliasing and not about acquisition
        order, so a \texttt{Mutex} guard whose lifetime the borrow checker tracks perfectly says
        nothing about whether two threads take their locks in a compatible
        order~\cite{zhang2026deadlock}. The lost update escapes ThreadSanitizer for the same
        reason it escapes the borrow checker: where every access is correctly synchronised,
        there is no unsynchronised access to find. So adopting Rust retires one category of
        concurrency defect and leaves untouched the category that fails silently, and a
        migration plan reading ``memory-safe'' as ``concurrency-safe'' would mis-budget its
        review effort. One case reaches past the ten injected programs and narrows the credit
        due on both sides. A specification gap let an operation be submitted outside the
        service's running window, and it survived every test, every parity vector and all three
        sanitisers because nothing ever called that path, which is the weakness of detection
        conditional on coverage. What turned undefined behaviour into a diagnosable failure on
        the Rust side was an ordinary bounds check and not a proof, so the property worth
        budgeting for is ``checked by default'' and not ``ruled out at compile time''. The root
        cause belongs to the method: a prohibition stated without a mechanism behind it leaves
        each implementation to fail in its own idiom.
  \item RQ3 splits by resource and the splits do not point the same way, while RQ4 comes out
        cheap in the machine and expensive in the people. The memory penalty reproduces under
        both allocators and is credible beyond its $p$-value, peak resident set being a
        run-level maximum insensitive to the jitter that destroys a p99, and it is a few MiB of
        fixed cost against an allocator effect an order of magnitude larger. One close study
        reports the opposite sign, with Rust lowest and the gap widening as concurrency
        rises~\cite{karlback2021ran}. The two are compatible: their C++ runs on Boost.Fibers,
        which allocates a stack per fiber, and their Rust on Tokio, where a task is a small
        heap-allocated future, so what they measured is the per-context cost of two concurrency
        runtimes. Two further studies agree with the direction found
        here~\cite{olsson2023ranmgmt,sundvisson2026drivers}. On migration the dominant cost is
        human~\cite{fulton2021adopting,crichton2024learning}, and confining \texttt{unsafe} to
        the seam bounds what has to be reviewed and not what can be corrupted, since host and
        component share one address space~\cite{rivera2021galeed,kirth2022pkrusafe}. The build and
        footprint axes point the same way as the seam: an organisation worried about either should
        audit its dependency budget before blaming the compiler, and the one axis where C++ leads
        outright is the edit-compile loop for changes that stay inside a translation unit.
  \item The three hypotheses were committed before any measurement, and stating their
        disposition is what separates a test from a description. H1 is supported below
        saturation and refuted above it, and not in the direction its wording anticipated, since
        the difference appeared in overload and not in the far tail H1 actually named. It was
        also not well posed, because ``performance-equivalent'' reads as one quantity and names
        at least three. H2 holds without exception for the five injected classes of its kind and
        fails in the looser reading its wording invites, so it is accepted for memory safety and
        data-race freedom and rejected for concurrency correctness at large. H3 is supported by
        the bound on a cost the measurement could have missed and not by the failure to detect
        one. Three threats stay load-bearing behind all of it. One author wrote both
        implementations, and the parity suite bounds behaviour and not
        performance~\cite{emberg2023asn1}. The end-to-end instrument measures the design
        faithfully and the language barely at all. And the overload divergence is replicated
        with its mechanism narrowed and not closed.
  \item The guidelines are the practical form of all this, and each carries the limit of the
        finding behind it. Choose a first component for the narrowness and the scalarity of its
        interface instead of for its size, because the redesign of a rich interface is the
        schedule risk and the boundary itself is not. Keep the C ABI in its own crate and the
        core free of \texttt{unsafe}, which makes the reviewable surface a countable
        thing~\cite{evans2020rust}. Leave the concurrency review budget where it is. Prefer a
        shared-nothing structure in either language, since the strongest concurrency result here
        belongs to the design and to neither compiler. Plan the learning curve as a training
        cost and instrument it, so the second component can be scheduled on measured effort.
        Measure the component as well as the round-trip, and sweep past the point where the
        service stops coping, which is the guideline this study would most like to have followed
        from the start.
\end{itemize}

\chapter{Conclusion and Future Work\label{chapter:conclusions}}

\begin{itemize}
  \item The thesis compared modern C++ and Rust on a UE Context Management service modelled on a
        real gNB-CU-CP component, along four dimensions: tail-latency predictability,
        compile-time defect prevention, safety-model overhead, and incremental-adoption cost.
        Both implementations were built against one frozen specification and driven along
        identical code paths, which is what lets a language effect be told apart from a
        difference of design or of skill.
  \item On safety the difference is one of kind and it has a definite edge, and the edge stops
        where the catalogue says it does. Safe Rust rejects all five injected memory-safety and
        data-race classes at compile time, turning hazards that C++ can only hope to catch with
        instrumented tests into programs that cannot be written. All five boundary classes
        compile and misbehave in safe Rust as they do in C++, and on two of them the C++ tooling
        reports what Rust's does not. Stated precisely, Rust eliminates data races and aliasing
        errors. It does not eliminate race conditions, lock-ordering errors, leaks, wrapping
        arithmetic or the blast radius of a failed worker.
  \item On performance the evidence divides twice before it divides by resource, and getting
        both divisions right was this study's most consequential correction to itself. It
        divides first by what is being measured, because an end-to-end operation here is
        $\BudgetCoreShareMax\%$ language and the rest an operating-system handoff, so an
        end-to-end null bounds nothing about either compiler, while the same operation measured
        apart is $\BudgetCoreRatio\times$ faster in \BudgetCoreFaster{}. It divides second by
        whether the service is coping. Below saturation the two are equivalent at the median and
        their far tail is unresolvable on a virtualised host, and past saturation
        \OverloadSustainedFaster{} sustains $\OverloadThrRatio\times$ the throughput of
        \OverloadSustainedSlower{} while drawing $\OverloadCpuPerOpRatio\times$ the CPU per
        completed operation. On memory the Rust build carries a few MiB more fixed footprint
        under both allocators, an order of magnitude less than the allocator choice costs.
  \item What follows is a narrow claim, narrower than this thesis expected to make. An
        unmodified Rust core can be driven from C++ across a thin seam whose cost on a real
        operation is bounded at $\FfiOpBoundPct\%$, with all \texttt{unsafe} confined to the seam
        and the core verified free of it. For stateful, concurrent, partitionable control-plane
        services of this kind, and at nominal load, the choice between the two languages is not a
        performance decision. It is a decision about which classes of defect an organisation
        wants to make unwriteable, and about whether it will pay the engineering and learning
        cost of getting them. Under overload it becomes a performance decision again, and one
        that currently favours C++. What the thesis adds is therefore an instrument and not a
        verdict: a frozen specification, a parity-tested pair of implementations, a measurement
        framework that reaches past the saturation knee, and a defect catalogue that maps where
        the guarantee stops.
  \item Future work answers the limitations one at a time. The immediate priority is the
        bare-metal campaign, whose target is now narrow: the far tail below saturation, with
        hardware counters and allocation profiles collected on the same visit, which is what
        would turn the memory result from an observation into an explanation. Beyond it,
        replaying anonymised production traces would validate the operation mix, a subsystem
        with genuinely shared mutable state would exercise the two languages differently than a
        shared-nothing design does, evaluating Rust async runtimes against the thread-per-shard
        model would separate the runtime from the language in the tail
        result~\cite{arvidsson2025async}, and measuring what in-process isolation of the two
        heaps costs would price what protecting an \texttt{unsafe}-free core is
        worth~\cite{rivera2021galeed}. A longer case study tracking defect density and review
        effort during an incremental adoption is what would settle the human cost this study can
        only name.
\end{itemize}

% Declaration of generative-AI use, required by the University when a language model has
% been used. Shared with Ha_Vien_Mastersthesis_2026.tex: edit ai_statement.tex, not a copy.
%% Declaration of generative-AI use. The University of Helsinki requires a student who
%% has used a language model in preparing a thesis to state in writing which model was
%% used and in what way; the model may not be named as an author or cited as a source.
%%
%% Shared verbatim by Thesis_POC.tex and Thesis_v1.tex so the two cannot drift apart.
%% It uses no generated macros and no chapter numbers, because the two documents number
%% their chapters differently.
%%
%% The order of the paragraphs follows the University's basic reporting instruction:
%% what the use covered, what it was for, which tool, how it performed and how its
%% output was checked, and what the author took from it.

\cleardoublepage
\phantomsection
\addcontentsline{toc}{chapter}{Statement on the Use of AI Tools}
\chapter*{Statement on the Use of AI Tools}
\markboth{Statement on the Use of AI Tools}{Statement on the Use of AI Tools}

\noindent\textbf{Tool.} Anthropic's Claude Opus~5, used through the Cursor development
environment, is the only generative AI system used in the preparation of this thesis. It was
used throughout the work, from June to August~2026. Following University of Helsinki policy the
model is neither named as an author nor cited as a source: it cannot take responsibility for
what it produces, and it cannot reliably vouch for the literature it paraphrases.

\noindent\textbf{What the assistance covered.} It was extensive, and is declared as such. The
model drafted and revised the running prose of this thesis and proposed the order in which its
arguments are made. It wrote and refactored the two implementations of the UE Context
Management service, the load generator, the fault-injection catalogue, and the interoperability
prototype. It wrote the analysis code behind the bootstrap confidence intervals, the exact
paired permutation tests, and the Holm correction, together with the figure-generating and
build scripts and the \LaTeX{} sources. It also suggested candidate literature.
% ==================== ACTION REQUIRED BEFORE SUBMISSION ==========================
% The abstract of Thesis_POC.tex was drafted with model assistance. The University's
% thesis guide states that a language model may not be used when preparing the thesis
% summary, and in this programme the summary also serves as the maturity test, where
% AI use is never permitted. Rewrite the abstract unaided before submitting; until
% that is done the closing sentence of this paragraph is not yet true of Thesis_POC.
% =================================================================================
The abstract was written by the author unaided, as the University's thesis guide requires of
the thesis summary.

\noindent\textbf{What the model did not settle.} The research problem, the four research
questions, and the three hypotheses are the author's, as is the control on which the study's
value rests: freezing one language-neutral specification and one concurrency model, so that the
language remains the only systematic difference between the two implementations. So is the
decision that keeps RQ2 falsifiable, namely to carry two defect classes -- a lock-order deadlock
and a lost update -- that lie outside Rust's compile-time guarantee and can therefore embarrass
the premise rather than confirm it. Every measurement reported here was run by the author on the
hosts documented in this thesis; no measured value came from a language model. Each measured
number reaches the text through a macro regenerated from the raw logs, so a figure cannot be
transcribed by hand, by the author or by the model. The readings of those numbers are the
author's as well, including the two restraints the evidence imposes: that a far tail a
virtualised host cannot resolve must not be reported as equivalence, and that a sweep dominated
by cross-thread handoff cannot bound an effect of the language.

\noindent\textbf{How the output was checked.} Model output was treated as a draft to be
falsified rather than as a result. Claims were checked against the sources cited for them: every
bibliography entry was resolved against Crossref, DataCite, or the issuing body, and each was
traced to the page or clause carrying the claim made of it. That audit found four passages where
the prose overstated or misplaced its source, one of them a sentence inside quotation marks that
does not appear in the work it was attributed to; all four were corrected, and the audit is kept
with the artefacts so that the correction remains reviewable. Code was checked by running it:
the two implementations are held to one conformance suite whose transcripts are diffed line by
line, and the fault catalogue fixes the expected outcome of each defect before it is injected,
so that a surprise appears as a discrepancy instead of being explained away afterwards. Nothing
is asserted here on the model's authority.

\noindent\textbf{What the author took from it.} The model is dependable at mechanical work --
transposing an idiom between two languages, restructuring a table, keeping a build honest -- and
undependable in a specific and predictable way: it will supply a plausible citation for a claim,
a plausible cause for a measurement, and a confident sentence where the data supports a hedge.
Its fluency is not evidence. The lasting lesson of the work was accordingly less about either
language than about verification. The effort saved in drafting is spent again, and rightly, in
checking, and the parts of this thesis that are defensible are the parts that were checked
against a source, a specification, or a measurement.

\noindent\textbf{Responsibility.} The author is responsible for the whole of this thesis: its
language, its facts, its code, its measurements, and its conclusions. Where it is wrong, the
error is the author's.


% STEP 5:
% The bibliography style is selected with the biblatex options in the preamble.

\cleardoublepage %fixes the position of bibliography in bookmarks
\phantomsection

\addcontentsline{toc}{chapter}{\bibname} % This lines adds the bibliography to the ToC
\printbibliography

\end{document}
